\chapter{Análisis de requisitos}
\label{chap:requisitos}


\section{Descripción del sistema}
\label{sec:descripcion_sistema}

El sistema desarrollado es un generador automático de modelos de
características expresados en UVL, integrado en el módulo
\textit{generator} de UVLHub. El generador se organiza mediante un
modelo de configuración y una operación compatibles con Flamapy, y
permite configurar y ejecutar la generación desde la interfaz web de
la plataforma.

El objetivo principal del sistema es facilitar la creación automática de modelos de características expresados en UVL con distintos niveles de complejidad y expresividad. Estos modelos permiten disponer de conjuntos de prueba más completos para la evaluación, validación y mejora de herramientas relacionadas con la Ingeniería de Líneas de Productos Software (SPL), especialmente aquellas orientadas al análisis de variabilidad. De esta forma, el sistema aborda la necesidad de disponer de modelos con características avanzadas que permitan evaluar el comportamiento de dichas herramientas en diferentes escenarios.

La entrada principal del sistema consiste en una configuración definida por el usuario mediante los parámetros disponibles en la interfaz proporcionada por UVLHub. Estos parámetros permiten establecer las características del modelo que se desea generar, incluyendo aspectos relacionados con su estructura jerárquica, atributos y restricciones.

Como resultado, el sistema produce modelos expresados en el lenguaje UVL con los niveles de expresividad soportados. Estos modelos pueden visualizarse directamente tras finalizar el proceso de generación dentro de la plataforma UVLHub y también pueden descargarse como ficheros independientes con extensión \texttt{.uvl} para su posterior utilización en otras herramientas del ecosistema de variabilidad software.

El sistema puede dividirse principalmente en los siguientes subsistemas:

\begin{itemize}

    \item \textbf{Subsistema de configuración del generador:} responsable de gestionar los parámetros definidos por el usuario mediante la interfaz web proporcionada en UVLHub. Este subsistema permite establecer las características del modelo que se desea generar y preparar la configuración necesaria para ejecutar el proceso desde el navegador.

    \item \textbf{Subsistema de generación automática:} constituye el núcleo funcional del sistema y es el encargado de producir los modelos de características a partir de la configuración proporcionada, haciendo uso de los componentes y mecanismos disponibles dentro del \textit{framework} Flamapy.

    \item \textbf{Subsistema de integración con UVLHub:} encargado de conectar la operación generadora desarrollada con el módulo \textit{generator} de UVLHub, permitiendo que la funcionalidad de la herramienta pueda ser utilizada dentro de la plataforma mediante los mecanismos definidos por esta.

\end{itemize}

El actor principal es el usuario investigador o desarrollador, que
configura el generador y obtiene modelos UVL para experimentación,
evaluación o validación. El acceso a esta funcionalidad no requiere
autenticación previa.

Las herramientas externas de análisis pueden utilizar los ficheros
descargados, pero su uso posterior no constituye una interacción
directa con el generador descrito en estos casos de uso.

\section{Requisitos de información}
\label{sec:ri}



El sistema gestiona la información necesaria para configurar el proceso de generación automática de modelos de características. Esta información se representa mediante un conjunto de entidades de configuración que permiten definir los parámetros utilizados durante la generación.

La entidad principal del modelo conceptual es \texttt{FmgeneratorModel}, que agrupa las diferentes configuraciones específicas del generador. Cada una de estas configuraciones representa un aspecto concreto del modelo que se desea generar, como la nomenclatura, los niveles de expresividad UVL soportados, las características, la estructura jerárquica, las restricciones o los atributos.

Estas entidades son utilizadas posteriormente por la operación de generación automática para producir modelos de características expresados en UVL. La operación de generación no se considera una entidad de información, sino una funcionalidad del sistema que será descrita en los requisitos funcionales.

\subsubsection*{RI-01: \texttt{FmgeneratorModel}}

La entidad \texttt{FmgeneratorModel} representa la configuración global del generador automático de modelos de características. Contiene los parámetros generales del proceso y las configuraciones específicas asociadas a cada aspecto del modelo generado.

La Tabla~\ref{tab:ri-fmgenerator-model} recoge sus parámetros generales y las configuraciones que agrupa.

\begin{longtable}{@{}P{5.3cm} P{\dimexpr\textwidth-5.3cm-2\tabcolsep\relax}@{}}
\toprule
\textbf{Atributo y tipo} & \textbf{Descripción} \\
\midrule
\endhead


\path{num_models}\newline{\footnotesize\texttt{int}} & Número de modelos que se desean generar. \\

\midrule
\path{seed}\newline{\footnotesize\texttt{int}} & Semilla utilizada para controlar la generación aleatoria. \\

\midrule
\path{ensure_satisfiable}\newline{\footnotesize\texttt{bool}} &
Indica si debe aplicarse la generación booleana con comprobación de
satisfacibilidad mediante SAT. \\

\midrule
\path{naming}\newline{\footnotesize\texttt{NamingConfig}} & Configuración relacionada con la nomenclatura de los elementos generados. \\

\midrule
\path{levels}\newline{\footnotesize\texttt{LevelConfig}} & Configuración de los niveles de expresividad soportados por UVL. \\

\midrule
\path{features}\newline{\footnotesize\texttt{FeaturesConfig}} & Configuración relacionada con la generación de características. \\

\midrule
\path{hierarchy}\newline{\footnotesize\texttt{HierarchyConfig}} & Configuración de la estructura jerárquica del modelo. \\

\midrule
\path{constraints}\newline{\footnotesize\texttt{ConstraintsConfig}} & Configuración de las restricciones generadas. \\

\midrule
\path{attributes}\newline{\footnotesize\texttt{AttributesConfig}} & Configuración de los atributos asociados a las características. \\

\bottomrule
\caption{Atributos del requisito de información RI-01.}
\label{tab:ri-fmgenerator-model}
\end{longtable}


\subsubsection*{RI-02: \texttt{NamingConfig}}

La entidad \texttt{NamingConfig} representa la configuración utilizada para definir la nomenclatura de los elementos generados por el sistema.

La Tabla~\ref{tab:ri-naming-config} presenta sus parámetros de nomenclatura.

\begin{longtable}{@{}P{5.3cm} P{\dimexpr\textwidth-5.3cm-2\tabcolsep\relax}@{}}
\toprule
\textbf{Atributo y tipo} & \textbf{Descripción} \\
\midrule
\endhead

\path{name_prefix}\newline{\footnotesize\texttt{str}} & Prefijo utilizado para los nombres generados. \\

\midrule
\path{include_feature_count_suffix}\newline{\footnotesize\texttt{bool}} & Indica si se incluye el número de características en el nombre del modelo generado. \\

\midrule
\path{include_constraint_count_suffix}\newline{\footnotesize\texttt{bool}} & Indica si se incluye el número de restricciones en el nombre del modelo generado. \\

\bottomrule
\caption{Atributos del requisito de información RI-02.}
\label{tab:ri-naming-config}
\end{longtable}


\subsubsection*{RI-03: \texttt{LevelConfig}}

La entidad \texttt{LevelConfig} representa los niveles de expresividad UVL que estarán disponibles durante el proceso de generación.

La Tabla~\ref{tab:ri-level-config} recoge los niveles configurables.

\begin{longtable}{@{}P{5.3cm} P{\dimexpr\textwidth-5.3cm-2\tabcolsep\relax}@{}}
\toprule
\textbf{Atributo y tipo} & \textbf{Descripción} \\
\midrule
\endhead

\path{boolean_level}\newline{\footnotesize\texttt{bool}} &
Representa el nivel booleano, que constituye la base obligatoria
de los modelos generados. \\

\midrule
\path{arithmetic_level}\newline{\footnotesize\texttt{bool}} & Indica si se habilitan elementos aritméticos en los modelos generados. \\

\midrule
\path{type_level}\newline{\footnotesize\texttt{bool}} &
Indica si se permite generar características tipadas y las
construcciones de UVL asociadas al nivel de tipos. \\

\midrule
\path{group_cardinality}\newline{\footnotesize\texttt{bool}} & Indica si se habilitan cardinalidades de grupo. \\

\midrule
\path{feature_cardinality}\newline{\footnotesize\texttt{bool}} & Indica si se habilitan cardinalidades individuales de características. \\

\midrule
\path{aggregate_functions}\newline{\footnotesize\texttt{bool}} & Indica si se permite la generación de funciones agregadas. \\

\midrule
\path{string_constraints}\newline{\footnotesize\texttt{bool}} & Indica si se habilitan restricciones sobre valores de tipo cadena. \\

\bottomrule
\caption{Atributos del requisito de información RI-03.}
\label{tab:ri-level-config}
\end{longtable}


\subsubsection*{RI-04: \textbf{FeaturesConfig}}

La entidad \textbf{FeaturesConfig} contiene los parámetros relacionados con la generación de características del modelo.

La Tabla~\ref{tab:ri-features-config} presenta los límites y distribuciones aplicables a las características.

\begin{longtable}{@{}P{5.3cm} P{\dimexpr\textwidth-5.3cm-2\tabcolsep\relax}@{}}
\toprule
\textbf{Atributo y tipo} & \textbf{Descripción} \\
\midrule
\endhead

\path{min_features}\newline{\footnotesize\texttt{int}} &
Número mínimo de características adicionales a la raíz que se
solicita generar. \\

\midrule
\path{max_features}\newline{\footnotesize\texttt{int}} &
Número máximo de características adicionales a la raíz que se
solicita generar. \\

\midrule
\path{dist_boolean}\newline{\footnotesize\texttt{float}} & Probabilidad de generación de características booleanas. \\

\midrule
\path{dist_integer}\newline{\footnotesize\texttt{float}} & Probabilidad de generación de características enteras. \\

\midrule
\path{dist_real}\newline{\footnotesize\texttt{float}} & Probabilidad de generación de características reales. \\

\midrule
\path{dist_string}\newline{\footnotesize\texttt{float}} & Probabilidad de generación de características de tipo cadena. \\

\midrule
\path{min_feature_cardinality}\newline{\footnotesize\texttt{int / list[int]}} & Valor mínimo permitido para las cardinalidades de características. \\

\midrule
\path{max_feature_cardinality}\newline{\footnotesize\texttt{int / list[int]}} & Valor máximo permitido para las cardinalidades de características. \\

\midrule
\path{prob_feature_cardinality}\newline{\footnotesize\texttt{float}} & Probabilidad de aplicar cardinalidades a las características. \\

\bottomrule
\caption{Atributos del requisito de información RI-04.}
\label{tab:ri-features-config}
\end{longtable}


\subsubsection*{RI-05: \texttt{HierarchyConfig}}

La entidad \texttt{HierarchyConfig} representa la configuración relacionada con la estructura jerárquica del modelo generado.

La Tabla~\ref{tab:ri-hierarchy-config} recoge los parámetros de profundidad y de relaciones jerárquicas.

\begin{longtable}{@{}P{5.3cm} P{\dimexpr\textwidth-5.3cm-2\tabcolsep\relax}@{}}
\toprule
\textbf{Atributo y tipo} & \textbf{Descripción} \\
\midrule
\endhead

\path{max_tree_depth}\newline{\footnotesize\texttt{int}} & Profundidad máxima permitida para la jerarquía del modelo. \\

\midrule
\path{dist_optional}\newline{\footnotesize\texttt{float}} & Probabilidad de generar relaciones opcionales. \\

\midrule
\path{dist_mandatory}\newline{\footnotesize\texttt{float}} & Probabilidad de generar relaciones obligatorias. \\

\midrule
\path{dist_alternative}\newline{\footnotesize\texttt{float}} & Probabilidad de generar relaciones alternativas. \\

\midrule
\path{dist_or}\newline{\footnotesize\texttt{float}} & Probabilidad de generar relaciones OR. \\

\midrule
\path{group_cardinality_min}\newline{\footnotesize\texttt{int}} & Cardinalidad mínima para grupos de características. \\

\midrule
\path{group_cardinality_max}\newline{\footnotesize\texttt{int}} & Cardinalidad máxima para grupos de características. \\

\midrule
\path{dist_group_cardinality}\newline{\footnotesize\texttt{float}} & Probabilidad de generar relaciones con cardinalidad de grupo. \\

\bottomrule
\caption{Atributos del requisito de información RI-05.}
\label{tab:ri-hierarchy-config}
\end{longtable}


\subsubsection*{RI-06: \texttt{ConstraintsConfig}}

La entidad \texttt{ConstraintsConfig} contiene los parámetros utilizados para definir la generación de restricciones entre elementos del modelo.

La Tabla~\ref{tab:ri-constraints-config} presenta los límites y distribuciones empleados para generar restricciones.

\begin{longtable}{@{}P{5.3cm} P{\dimexpr\textwidth-5.3cm-2\tabcolsep\relax}@{}}
\toprule
\textbf{Atributo y tipo} & \textbf{Descripción} \\
\midrule
\endhead

\path{min_constraints}\newline{\footnotesize\texttt{int}} & Número mínimo de restricciones generadas. \\

\midrule
\path{max_constraints}\newline{\footnotesize\texttt{int}} & Número máximo de restricciones generadas. \\

\midrule
\path{extra_constraint_representativeness}\newline{\footnotesize\texttt{int}} & Factor utilizado para aumentar la representatividad de las restricciones. \\

\midrule
\path{min_vars_per_constraint}\newline{\footnotesize\texttt{int}} & Número mínimo de variables utilizadas en una restricción. \\

\midrule
\path{max_vars_per_constraint}\newline{\footnotesize\texttt{int}} & Número máximo de variables utilizadas en una restricción. \\

\midrule
\path{ctc_dist_boolean}\newline{\footnotesize\texttt{float}} & Probabilidad de generar restricciones booleanas. \\

\midrule
\path{ctc_dist_integer}\newline{\footnotesize\texttt{float}} & Probabilidad de generar restricciones enteras. \\

\midrule
\path{ctc_dist_real}\newline{\footnotesize\texttt{float}} & Probabilidad de generar restricciones reales. \\

\midrule
\path{ctc_dist_string}\newline{\footnotesize\texttt{float}} & Probabilidad de generar restricciones sobre cadenas. \\

\midrule
\path{prob_not}\newline{\footnotesize\texttt{float}} & Probabilidad de aplicar negación en restricciones booleanas. \\

\midrule
\path{prob_and}\newline{\footnotesize\texttt{float}} & Probabilidad de utilizar operadores AND. \\

\midrule
\path{prob_or_ct}\newline{\footnotesize\texttt{float}} & Probabilidad de utilizar operadores OR. \\

\midrule
\path{prob_implication}\newline{\footnotesize\texttt{float}} & Probabilidad de utilizar implicaciones. \\

\midrule
\path{prob_equivalence}\newline{\footnotesize\texttt{float}} & Probabilidad de utilizar equivalencias. \\

\midrule
\path{prob_sum}\newline{\footnotesize\texttt{float}} & Probabilidad de utilizar operaciones de suma. \\

\midrule
\path{prob_substract}\newline{\footnotesize\texttt{float}} & Probabilidad de utilizar operaciones de resta. \\

\midrule
\path{prob_multiply}\newline{\footnotesize\texttt{float}} & Probabilidad de utilizar operaciones de multiplicación. \\

\midrule
\path{prob_divide}\newline{\footnotesize\texttt{float}} & Probabilidad de utilizar operaciones de división. \\

\midrule
\path{prob_equals}\newline{\footnotesize\texttt{float}} & Probabilidad de utilizar comparaciones de igualdad. \\

\midrule
\path{prob_less}\newline{\footnotesize\texttt{float}} & Probabilidad de utilizar comparaciones menor que. \\

\midrule
\path{prob_greater}\newline{\footnotesize\texttt{float}} & Probabilidad de utilizar comparaciones mayor que. \\

\midrule
\path{prob_less_equals}\newline{\footnotesize\texttt{float}} & Probabilidad de utilizar comparaciones menor o igual que. \\

\midrule
\path{prob_greater_equals}\newline{\footnotesize\texttt{float}} & Probabilidad de utilizar comparaciones mayor o igual que. \\

\midrule
\path{prob_sum_function}\newline{\footnotesize\texttt{float}} & Probabilidad de utilizar funciones agregadas de suma. \\

\midrule
\path{prob_avg_function}\newline{\footnotesize\texttt{float}} & Probabilidad de utilizar funciones agregadas de promedio. \\

\midrule
\path{prob_len_function}\newline{\footnotesize\texttt{float}} & Probabilidad de utilizar funciones de longitud sobre cadenas. \\

\bottomrule
\caption{Atributos del requisito de información RI-06.}
\label{tab:ri-constraints-config}
\end{longtable}


\subsubsection*{RI-07: \texttt{AttributesConfig}}

La entidad \texttt{AttributesConfig} representa la configuración relacionada con los atributos asociados a las características generadas.

La Tabla~\ref{tab:ri-attributes-config} recoge los parámetros de generación automática y manual de atributos.

\begin{longtable}{@{}P{5.3cm} P{\dimexpr\textwidth-5.3cm-2\tabcolsep\relax}@{}}
\toprule
\textbf{Atributo y tipo} & \textbf{Descripción} \\
\midrule
\endhead

\path{random_attributes}\newline{\footnotesize\texttt{bool}} & Indica si los atributos deben generarse automáticamente o definirse manualmente. \\

\midrule
\path{min_attributes}\newline{\footnotesize\texttt{int / None}} & Número mínimo de atributos generados en modo aleatorio. \\

\midrule
\path{max_attributes}\newline{\footnotesize\texttt{int / None}} & Número máximo de atributos generados en modo aleatorio. \\

\midrule
\path{dist_boolean_atr}\newline{\footnotesize\texttt{float}} & Probabilidad de seleccionar el tipo booleano en la generación
aleatoria de atributos. \\

\midrule
\path{dist_integer_atr}\newline{\footnotesize\texttt{float}} & Probabilidad de seleccionar el tipo entero en la generación
aleatoria de atributos. \\

\midrule
\path{dist_real_atr}\newline{\footnotesize\texttt{float}} & Probabilidad de seleccionar el tipo real en la generación
aleatoria de atributos. \\

\midrule
\path{dist_string_atr}\newline{\footnotesize\texttt{float}} & Probabilidad de seleccionar el tipo cadena en la generación
aleatoria de atributos. \\

\midrule
\path{attributes_list}\newline{\footnotesize\texttt{list[dict]}} & Lista de atributos definidos manualmente para la generación. Cada
definición puede incluir su tipo, dominio, probabilidad de asociación y
participación en restricciones. \\

\midrule
\path{attribute_attach_probs}\newline{\footnotesize\texttt{list[float]}} & Probabilidades de asociación de atributos a características. \\

\midrule
\path{attribute_in_constraints}\newline{\footnotesize\texttt{list[bool]}} & Indica si cada atributo definido manualmente puede participar en
restricciones. \\

\bottomrule
\caption{Atributos del requisito de información RI-07.}
\label{tab:ri-attributes-config}
\end{longtable}

\section{Requisitos funcionales}
\label{sec:rf}


El sistema debe proporcionar las funcionalidades necesarias para permitir la configuración, ejecución y obtención de modelos de características generados automáticamente. Los requisitos funcionales identificados se muestran en la Tabla~\ref{tab:requisitos-funcionales}.

\begin{longtable}{@{}P{1.25cm} P{3.2cm}
P{\dimexpr\textwidth-5.7cm-6\tabcolsep\relax} P{1.25cm}@{}}

\toprule
\textbf{ID} & \textbf{Nombre} & \textbf{Descripción} & \textbf{Prior.} \\
\midrule
\endhead

RF-01 &
Configuración del generador &
El sistema debe proporcionar un asistente de seis pasos para configurar
la generación. El usuario debe poder indicar: el número de modelos, entre 1 y
1000; una semilla entera positiva; un prefijo opcional; los \textit{major
levels} y \textit{minor levels} de UVL; los límites y distribuciones del
\textit{feature tree}; el número, composición y operadores de las
\textit{cross-tree constraints}; la generación automática o manual de
\textit{attributes}; y las opciones de satisfacibilidad y nomenclatura de los
ficheros. Los parámetros y sus dependencias se detallan en las
reglas RN-01 a RN-06.
&
Alta \\[0.8em]

RF-02 &
Validación de la configuración &
El sistema debe validar cada paso antes de permitir el avance. Debe comprobar
que los valores enteros sean válidos, que los límites mínimos no superen los
máximos, que las probabilidades pertenezcan al intervalo $[0,1]$, que las
distribuciones activas sumen 1 y que los \textit{minor levels} sean compatibles
con sus \textit{major levels}. Cuando una validación falle, el sistema debe
mantener los valores introducidos, permanecer en el mismo paso y mostrar el
mensaje asociado al campo incorrecto. Por ejemplo:
\texttt{Number of models cannot exceed 1000.},
\texttt{Requires Arithmetic level.} o
\texttt{Current sum: 0.8000. Must be 1.0.}
&
Alta \\[0.8em]

RF-03 &
Generación automática de modelos UVL &
En el flujo ordinario, el sistema debe transformar la configuración
validada en una instancia de \texttt{FmgeneratorModel} en el navegador
y ejecutar \texttt{GenerateFeatureModel.execute()} para cada modelo
solicitado. La operación debe producir un \texttt{FeatureModel} y aplicar la semilla
configurada. Posteriormente, el resultado debe serializarse mediante
\texttt{UVLWriter}. Cuando se active la comprobación de
satisfacibilidad, la generación y la verificación mediante SAT deben
realizarse en el servidor. La semilla efectiva debe depender de la
semilla base, el índice del modelo y, cuando proceda, el índice del
intento. Los errores deben comunicarse al usuario sin presentar un
modelo incompleto como resultado válido.
&
Alta \\[0.8em]

RF-04 &
Soporte de niveles de expresividad UVL &
El sistema debe generar modelos del \textit{Boolean level} y permitir la
activación del \textit{Arithmetic level}, el \textit{Type level} y los
\textit{minor levels} de \textit{group cardinality}, \textit{feature
cardinality}, \textit{aggregate functions} y \textit{string constraints}. Cuando la comprobación de satisfacibilidad esté desactivada, las
opciones seleccionadas deben habilitar las construcciones correspondientes
y reflejarse en las declaraciones \texttt{include} pertinentes. Su
activación no garantiza que cada construcción aparezca en todos los
modelos generados:
\texttt{Boolean.group-cardinality},
\texttt{Arithmetic.feature-cardinality},
\texttt{Arithmetic.aggregate-function}, \texttt{Arithmetic.*} o
\texttt{Type.string-constraints}. Cuando esté activada, se aplicará
la limitación al nivel booleano descrita en RN-06.
&
Alta \\[0.8em]

RF-05 &
Visualización del modelo generado &
Tras completar correctamente la generación, el sistema debe mostrar
en UVLHub los resultados obtenidos, incluido el nombre del fichero
y los recuentos de características y restricciones. El usuario debe
poder abrir la vista del contenido UVL de cada modelo. Estas acciones
solo deben estar disponibles para modelos generados correctamente;
si la generación falla, el sistema debe informar del error.
&
Media \\[0.8em]

RF-06 &
Descarga de los modelos generados &
El sistema debe permitir descargar los resultados como ficheros con extensión
\texttt{.uvl}. Cuando se generen varios modelos, estos deben incluirse en un
archivo comprimido. El nombre de cada fichero debe utilizar el prefijo indicado
o \texttt{fm} como valor predeterminado y aplicar, cuando hayan sido
seleccionados, los sufijos con el número de características y restricciones. Si
no se emplean sufijos y se genera más de un modelo, el sistema debe añadir su
índice para evitar nombres idénticos.
&
Alta \\

\bottomrule
\caption{Requisitos funcionales del sistema.}
\label{tab:requisitos-funcionales}
\end{longtable}

\section{Requisitos no funcionales}
\label{sec:rnf}

La Tabla~\ref{tab:requisitos-no-funcionales} reúne las condiciones de integración, compatibilidad, ejecución y calidad del sistema.

\begin{longtable}{@{}P{1.25cm} P{3.2cm}
P{\dimexpr\textwidth-5.7cm-6\tabcolsep\relax} P{1.25cm}@{}}

\toprule
\textbf{ID} & \textbf{Nombre} & \textbf{Descripción} & \textbf{Prior.} \\
\midrule
\endhead

RNF-01 &
Modularidad e integración &
El generador debe mantenerse como un componente independiente compatible con
la arquitectura de Flamapy. La operación
\texttt{GenerateFeatureModel} debe implementar la interfaz de operación de
Flamapy, poder localizarse mediante \texttt{DiscoverMetamodels} y utilizarse
desde UVLHub sin duplicar su implementación.
&
Alta \\[0.8em]

RNF-02 &
Compatibilidad sintáctica con UVL &
Los modelos serializados deben poder procesarse mediante
\texttt{UVLReader} sin errores sintácticos. Las declaraciones
\texttt{include} necesarias para las construcciones generadas deben conservarse
durante la serialización.
&
Alta \\[0.8em]

RNF-03 &
Ejecución en el navegador &
El flujo ordinario de generación debe ejecutarse en el navegador
mediante Pyodide y WebAssembly. Para ese flujo, el servidor almacena
la configuración como un diccionario en la sesión y la expone como
JSON (\textit{JavaScript Object Notation}); la instancia de
\texttt{FmgeneratorModel} se construye en el navegador. Cuando el
usuario activa la comprobación de satisfacibilidad, la generación
y la verificación mediante SAT se realizan en el servidor.
&
Alta \\[0.8em]

RNF-04 &
Usabilidad del asistente &
La configuración debe dividirse en seis pasos agrupados por funcionalidad. El
sistema debe conservar los valores al navegar hacia atrás o después de una
validación fallida, ocultar o desactivar las opciones que no sean aplicables y
mostrar los errores junto al parámetro que debe corregirse.
&
Media \\[0.8em]

RNF-05 &
Calidad del software &
La cobertura de sentencias ejecutables del conjunto de ficheros de
código de producción evaluados en la
Tabla~\ref{tab:cobertura_ficheros} debe ser igual o superior al
80\,\%. La solución debe superar la suite automatizada descrita en
el Capítulo~\ref{chap:pruebas}. Los workflows de integración
continua del repositorio upstream de UVLHub no forman parte de los
criterios de aceptación específicos de este TFG.
&
Alta \\

\bottomrule
\caption{Requisitos no funcionales del sistema.}
\label{tab:requisitos-no-funcionales}
\end{longtable}

\section{Reglas de negocio}
\label{sec:reglas_negocio}

La Tabla~\ref{tab:reglas-negocio} concreta las validaciones y dependencias aplicables a la configuración.

\begin{longtable}{@{}P{1.2cm} P{3.2cm}
P{\dimexpr\textwidth-4.4cm-4\tabcolsep\relax}@{}}

\toprule
\textbf{ID} & \textbf{Nombre} & \textbf{Descripción} \\
\midrule
\endhead

RN-01 &
Parámetros generales &
El número de modelos, \texttt{num\_models}, debe ser un entero comprendido
entre 1 y 1000. Los valores inferiores producen el mensaje
\texttt{Number of models must be at least 1.}; los superiores,
\texttt{Number of models cannot exceed 1000.}; y los valores no enteros,
\texttt{Number of models must be an integer.}

La semilla, \texttt{seed}, debe ser un entero positivo. En caso contrario se
muestra \texttt{Seed must be a positive integer.} El prefijo es opcional; si
está vacío, se normaliza al valor \texttt{fm}.
\\[1em]

RN-02 &
Dependencias entre niveles UVL &
El \textit{Boolean level} debe permanecer siempre activado. La selección del
\textit{Type level} activa automáticamente el \textit{Arithmetic level}. Los
\textit{minor levels} de \textit{feature cardinality} y
\textit{aggregate functions} requieren el \textit{Arithmetic level}; si este
no está activo, se muestra \texttt{Requires Arithmetic level.}

El nivel de \textit{string constraints} requiere el \textit{Type level}. Si se
selecciona sin él, se muestra \texttt{Requires Type level.} La
\textit{group cardinality} pertenece al \textit{Boolean level} y no requiere
activar niveles superiores.
\\[1em]

RN-03 &
Configuración del \textit{feature tree} &
Los parámetros \texttt{min\_features} y \texttt{max\_features} deben ser
enteros mayores o iguales que 1 y el mínimo no puede superar al máximo. Los
mensajes utilizados son \texttt{Min. features must be at least 1.},
\texttt{Max. features must be at least 1.} y
\texttt{Max. features must be >= Min. features.}

La profundidad \texttt{max\_tree\_depth} debe ser un entero comprendido entre
1 y \texttt{max\_features}; en caso contrario se muestra
\texttt{Maximum tree depth must be between 1 and Max. features.}

Las probabilidades de las relaciones \textit{optional}, \textit{mandatory},
\textit{alternative}, \textit{or} y, cuando esté habilitada,
\textit{group cardinality}, deben pertenecer a $[0,1]$ y sumar 1. Se admite
una tolerancia de 0,001 en la interfaz. Los errores muestran
\texttt{Value must be between 0 and 1.},
\texttt{The relation-distribution probabilities must sum to 1.0.} y el valor
actual mediante \texttt{Current sum: X. Must be 1.0.}

Cuando se active la \textit{group cardinality}, sus límites deben ser enteros
mayores o iguales que 0, el mínimo no puede superar al máximo y el máximo no
puede superar \texttt{max\_features}. Cuando se active la
\textit{feature cardinality}, sus límites también deben ser mayores o iguales
que 1, el mínimo no puede superar al máximo y
\texttt{prob\_feature\_cardinality} debe pertenecer a $[0,1]$.

La distribución de características booleanas, enteras, reales y de cadena debe
estar formada por probabilidades de $[0,1]$ cuya suma sea 1.
\\[1em]

RN-04 &
Configuración de restricciones &
Los parámetros \texttt{min\_constraints} y \texttt{max\_constraints} deben ser
enteros mayores o iguales que 1 y el mínimo no puede superar al máximo. Los
parámetros \texttt{min\_vars\_per\_constraint} y
\texttt{max\_vars\_per\_constraint} siguen la misma regla, y el máximo de
variables no puede superar \texttt{max\_features}.

\texttt{extra\_constraint\_representativeness} debe ser un entero mayor o igual
que 1 y menor o igual que \texttt{max\_vars\_per\_constraint}. Los mensajes
asociados son \texttt{Must be an integer.}, \texttt{Must be >= 1.} y
\texttt{Must be <= Max. vars per constraint.}

Las probabilidades de \texttt{and}, \texttt{or}, implicación y equivalencia
deben sumar 1. La probabilidad de negación se valida de forma independiente
dentro de $[0,1]$. Si el \textit{Arithmetic level} está activo, las
probabilidades de suma, resta, multiplicación y división deben sumar 1. Cuando
también estén activadas las \textit{aggregate functions}, se incorporan a esa
misma distribución las probabilidades de \texttt{sum()} y \texttt{avg()}.

Las probabilidades de los operadores de comparación deben sumar 1. La
distribución de tipos de restricciones debe sumar 1 considerando únicamente
los tipos disponibles según los niveles activos. Por último,
\texttt{prob\_len\_function} debe pertenecer a $[0,1]$ cuando estén activos el
\textit{Type level} y las \textit{string constraints}.

Ante una suma incorrecta, el sistema muestra el valor calculado y mensajes como
\texttt{The boolean-connective probabilities must sum to 1.0.},
\texttt{Arithmetic-level probabilities must sum to 1.0.},
\texttt{Comparison-operator probabilities must sum to 1.0.} o
\texttt{Active type probabilities must total 1.0.}
\\[1em]

RN-05 &
Configuración de \textit{attributes} &
En modo automático, \texttt{min\_attributes} y \texttt{max\_attributes} deben
ser enteros comprendidos entre 0 y 1000 y el mínimo no puede superar al máximo.
Los mensajes asociados son
\texttt{Min. attributes must be between 0 and 1000.},
\texttt{Max. attributes must be between 0 and 1000.} y
\texttt{Max. attributes must be >= Min.}

La distribución de \textit{attributes} booleanos, enteros, reales y de cadena
debe utilizar valores de $[0,1]$ y sumar 1 considerando únicamente los tipos
habilitados por los niveles seleccionados.

En modo manual, cada \textit{attribute} debe tener un nombre y una probabilidad
de asociación perteneciente a $[0,1]$. Si falta alguno de estos valores se
muestran \texttt{Attribute name is required.} o
\texttt{Attach probability is required.}

Un \textit{attribute} booleano debe permitir al menos uno de los valores
verdadero o falso. Los atributos enteros, reales y de cadena deben definir
límites mínimo y máximo numéricos, con el mínimo menor o igual que el máximo.
En los atributos de cadena, ambos límites deben ser no negativos.

Los atributos enteros o reales solo pueden participar en restricciones cuando
está activo el \textit{Arithmetic level}. Los atributos de cadena requieren el
\textit{Type level} y las \textit{string constraints}. En caso contrario se
muestran, respectivamente,
\texttt{Integer/Real attrs can be used in constraints only when Arithmetic
level is on.} y
\texttt{String attrs can be used in constraints only when Type level + String
constraints are on.}
\\[1em]

RN-06 &
Satisfacibilidad y nomenclatura de salida &
Cuando se active \texttt{ensure\_satisfiable}, la configuración
efectiva utilizada para generar los modelos se limitará al nivel
booleano. Se desestimarán los niveles aritmético y de tipos, sus
niveles menores y la cardinalidad de grupo, aunque se hubieran
seleccionado en el asistente. Tampoco se generarán atributos,
características tipadas ni restricciones ajenas al dominio booleano.

Cada modelo dispondrá de un máximo de 20 intentos. Cada candidato se
comprobará mediante SAT y solo se devolverá si se confirma que es
satisfacible. Si se agotan los intentos, se producirá
\texttt{SatisfiableModelGenerationError}, cuyo mensaje identifica el
número de modelo y los 20 intentos realizados. El error se comunicará
al usuario sin presentar el último candidato como resultado válido.

Cuando \texttt{ensure\_satisfiable} esté desactivado, se tendrán en
cuenta los niveles de expresividad, las cardinalidades y los atributos
habilitados en la configuración.

Los resultados se serializarán con extensión \texttt{.uvl}. El nombre
base será el prefijo configurado o \texttt{fm}. Los sufijos opcionales
tendrán la forma \texttt{\_Nf} y \texttt{\_Mc}, donde $N$ es el número
de características y $M$ el número de restricciones. Cuando no se
activen sufijos y \texttt{num\_models} sea mayor que 1, se añadirá el
índice del modelo.
\\

\bottomrule
\caption{Reglas de negocio del dominio.}
\label{tab:reglas-negocio}
\end{longtable}

\section{Restricciones del proyecto}
\label{sec:restricciones}

La Tabla~\ref{tab:res} recoge las restricciones tecnológicas, organizativas y temporales del proyecto.

\begin{longtable}{@{}P{1.4cm} P{3cm}
P{\dimexpr\textwidth-4.4cm-4\tabcolsep\relax}@{}}

\toprule
\textbf{ID} & \textbf{Tipo} & \textbf{Descripción} \\
\midrule
\endhead

RES-01 & Tecnológica &
El desarrollo debe integrarse dentro del ecosistema existente de herramientas de variabilidad software, utilizando como base el \textit{framework} Flamapy y la plataforma UVLHub. Esta condición limita la arquitectura de la solución, ya que el generador debe seguir los mecanismos de extensión definidos por Flamapy, adaptarse a la estructura proporcionada por UVLHub y utilizar las librerías propias del \textit{framework} y componentes externos necesarios para la generación y representación de modelos de características.
\\[0.8em]

RES-02 & Tecnológica / interoperabilidad &
Los modelos generados deben utilizar el lenguaje UVL como formato de salida, debido a que constituye el estándar empleado dentro del ecosistema de herramientas con el que se integra el proyecto.
\\[0.8em]

RES-03 & Organizativa &
El desarrollo debe realizarse dentro del contexto del grupo de investigación y de los proyectos software existentes asociados, respetando los procedimientos, recursos y estructuras de trabajo establecidos.
\\[0.8em]

RES-04 & Temporal &
El desarrollo y la entrega de la memoria deben ajustarse al calendario
académico del TFG, cuya fecha límite de entrega es el 7 de octubre
de 2026.
\\[0.8em]

\bottomrule

\caption{Restricciones del proyecto.}
\label{tab:res}
\end{longtable}

\section{Casos de uso}
\label{sec:casos_uso}

\subsection{Diagrama general de casos de uso}
\label{subsec:cu_diagrama}


El diagrama de casos de uso representa las principales interacciones entre los actores externos y el sistema generador de modelos UVL. El usuario investigador o desarrollador es el encargado de configurar los parámetros del generador y, una vez finalizado el proceso, puede visualizar los modelos generados dentro de la plataforma o descargarlos para su posterior utilización. El sistema se encarga de validar la configuración proporcionada y ejecutar el proceso de generación automática de modelos de características UVL. Los modelos descargados pueden ser utilizados posteriormente por herramientas externas de análisis de variabilidad para realizar procesos de evaluación y análisis sobre los modelos generados.

La Figura~\ref{fig:diagrama-casos-uso} muestra el diagrama general de casos de uso del sistema.

\begin{figure}[H]
    \centering
    \makebox[\textwidth][c]{
        \includegraphics[width=1\textwidth]{diagrama_casos_uso.pdf}
    }
    \caption{Diagrama general de casos de uso del sistema. Elaboración propia.}
    \label{fig:diagrama-casos-uso}
\end{figure}

\subsection{Especificación de casos de uso principales}
\label{subsec:cu_especificacion}

La Tabla~\ref{tab:uc01} especifica el caso de uso de configuración del generador.

\begin{table}[H]
\centering
\small
\renewcommand{\arraystretch}{1.35}
\begin{tabular}{@{}>{\bfseries\RaggedRight\arraybackslash}p{4cm} P{10cm}@{}}
\toprule
\multicolumn{2}{l}{\textbf{CU-01 --- Configurar generador}} \\
\midrule

Actores &
Usuario investigador o desarrollador. \\

Descripción &
Permite al usuario definir la configuración necesaria para el proceso de generación automática de modelos de características UVL. El usuario puede seleccionar los niveles de expresividad y establecer los parámetros relacionados con características, estructura jerárquica, restricciones y atributos. \\

Precondiciones &
El usuario debe acceder a la interfaz de configuración del generador disponible dentro de UVLHub. \\

Flujo principal &
1. El usuario accede al formulario de configuración del generador.

2. El usuario establece los parámetros asociados a la generación del modelo en cada paso del asistente.

3. El sistema registra la configuración proporcionada. \\

Flujo alternativo &
Si los valores introducidos no cumplen las restricciones establecidas, el sistema informa al usuario de los errores detectados durante la introducción de dichos valores o antes de pasar al siguiente paso. \\

Postcondiciones &
La configuración validada queda disponible para que el usuario
solicite la generación desde el último paso. \\

\bottomrule
\end{tabular}
\caption{Especificación del caso de uso CU-01.}
\label{tab:uc01}
\end{table}

La Tabla~\ref{tab:uc02} describe la validación que se realiza al avanzar por el asistente.

\begin{table}[H]
\centering
\small
\renewcommand{\arraystretch}{1.35}
\begin{tabular}{@{}>{\bfseries\RaggedRight\arraybackslash}p{4cm} P{10cm}@{}}
\toprule
\multicolumn{2}{l}{\textbf{CU-02 --- Validar configuración}} \\
\midrule

Actores &
Usuario investigador o desarrollador. \\

Descripción &
Permite comprobar que la configuración proporcionada por el usuario cumple las reglas de negocio definidas antes de ejecutar el proceso de generación automática. \\

Precondiciones &
Debe existir una configuración del generador proporcionada por el usuario. \\

Flujo principal &
1. El usuario solicita avanzar al siguiente paso del asistente.

2. El sistema comprueba los valores, las distribuciones y las
dependencias correspondientes al paso actual.

3. Si la configuración es válida, el sistema guarda los parámetros y
permite continuar. \\

Flujo alternativo &
Si alguna regla de negocio no se cumple, el sistema conserva los
valores introducidos, muestra los errores detectados e impide avanzar
hasta que se corrija la configuración. \\

Postcondiciones &
Si la configuración de cada paso es válida, se le permite al usuario pasar al siguiente paso. \\

\bottomrule
\end{tabular}
\caption{Especificación del caso de uso CU-02.}
\label{tab:uc02}
\end{table}

La Tabla~\ref{tab:uc03} especifica la solicitud y ejecución de la generación.

\begin{table}[H]
\centering
\small
\renewcommand{\arraystretch}{1.35}
\begin{tabular}{@{}>{\bfseries\RaggedRight\arraybackslash}p{4cm} P{10cm}@{}}
\toprule
\multicolumn{2}{l}{\textbf{CU-03 --- Generar modelo de características UVL}} \\
\midrule

Actores &
Usuario investigador o desarrollador. \\

Descripción &
Permite generar automáticamente un modelo de características expresado en UVL a partir de una configuración previamente validada. \\

Precondiciones &
La configuración proporcionada por el usuario debe haberse validado
en los pasos anteriores y el entorno necesario para el flujo de
generación seleccionado debe estar disponible. \\

Flujo principal &
1. Desde el último paso, el usuario solicita la generación y el
sistema inicia el flujo correspondiente a las opciones seleccionadas.

2. El sistema genera, para cada modelo y de forma aleatoria, la estructura jerárquica del modelo de características, incorporando los niveles de expresividad, atributos y restricciones pertinentes, cumpliendo con la configuración seleccionada por el usuario.

3. El sistema genera la representación final de los modelos en formato UVL y los almacena temporalmente para su visualización o descarga. \\

Flujo alternativo &
Si se produce un error, el sistema informa del fallo y no presenta
como válidos los modelos cuya generación o comprobación no se haya
completado. \\

Postcondiciones &
Existen uno o más modelos de características UVL generados correctamente disponibles para su posterior visualización o descarga. \\

\bottomrule
\end{tabular}
\caption{Especificación del caso de uso CU-03.}
\label{tab:uc03}
\end{table}

La Tabla~\ref{tab:uc04} describe la visualización de un resultado generado.

\begin{table}[H]
\centering
\small
\renewcommand{\arraystretch}{1.35}
\begin{tabular}{@{}>{\bfseries\RaggedRight\arraybackslash}p{4cm} P{10cm}@{}}
\toprule
\multicolumn{2}{l}{\textbf{CU-04 --- Visualizar modelo generado}} \\
\midrule

Actores &
Usuario investigador o desarrollador. \\

Descripción &
Permite al usuario consultar visualmente el modelo de características generado dentro de la plataforma UVLHub. \\

Precondiciones &
Debe existir al menos un modelo generado correctamente por el sistema. \\

Flujo principal &
1. El usuario selecciona la acción de visualizar uno de los
modelos generados.

2. El sistema abre una vista con su contenido UVL.

3. El usuario consulta el modelo y cierra la vista cuando termina. \\

Flujo alternativo &
Si no se ha podido realizar la generación previa a la visualización, el sistema lanza un error informativo. \\

Postcondiciones &
El usuario puede consultar el modelo generado dentro de la plataforma. \\

\bottomrule
\end{tabular}
\caption{Especificación del caso de uso CU-04.}
\label{tab:uc04}
\end{table}

La Tabla~\ref{tab:uc05} recoge las opciones de descarga de los resultados.

\begin{table}[H]
\centering
\small
\renewcommand{\arraystretch}{1.35}
\begin{tabular}{@{}>{\bfseries\RaggedRight\arraybackslash}p{4cm} P{10cm}@{}}
\toprule
\multicolumn{2}{l}{\textbf{CU-05 --- Descargar modelo generado}} \\
\midrule

Actores &
Usuario investigador o desarrollador. \\

Descripción &
Permite descargar individualmente los modelos como ficheros
\texttt{.uvl} o descargar el conjunto generado en un archivo
comprimido. \\

Precondiciones &
Debe existir al menos un modelo de características generado correctamente. \\

Flujo principal &
1. El usuario selecciona la descarga de un modelo concreto o la
descarga del conjunto de resultados.

2. El sistema prepara el fichero \texttt{.uvl} correspondiente o el
archivo comprimido con los modelos generados.

3. El navegador inicia la descarga del fichero preparado.
\\

Flujo alternativo &
Si el modelo no se encuentra disponible o se produce un error durante la generación del fichero, el sistema informa al usuario y cancela la descarga. \\

Postcondiciones &
El usuario dispone del modelo descargado en formato \texttt{.uvl}
o del archivo comprimido con el conjunto de modelos. \\

\bottomrule
\end{tabular}
\caption{Especificación del caso de uso CU-05.}
\label{tab:uc05}
\end{table}

\section{Matriz de trazabilidad}
\label{sec:trazabilidad}


\subsection{Requisitos funcionales \texorpdfstring{$\leftrightarrow$}{<->} casos de uso}
\label{subsec:trazabilidad_rf_cu}

La Tabla~\ref{tab:trazabilidad_rf_cu} relaciona cada requisito funcional con el caso de uso que lo desarrolla.

\begin{longtable}{P{1.7cm} *{5}{P{1.8cm}}}

\toprule
\textbf{} & \textbf{CU-01} & \textbf{CU-02} & \textbf{CU-03} & \textbf{CU-04} & \textbf{CU-05} \\
\midrule
\endhead

RF-01 & \texttimes & & & & \\

RF-02 & & \texttimes & & & \\

RF-03 & & & \texttimes & & \\

RF-04 & \texttimes & & \texttimes & & \\

RF-05 & & & & \texttimes & \\

RF-06 & & & & & \texttimes \\

\bottomrule

\caption{Matriz de trazabilidad entre requisitos funcionales y casos de uso.}
\label{tab:trazabilidad_rf_cu}

\end{longtable}

\subsection{Requisitos \texorpdfstring{$\leftrightarrow$}{<->} pruebas}
\label{subsec:trazabilidad_rf_pruebas}

La Tabla~\ref{tab:trazabilidad_rf_tests} recoge las pruebas y evidencias
utilizadas para verificar los requisitos funcionales y no funcionales.

\begingroup
\small
\setlength{\tabcolsep}{4pt}
\renewcommand{\arraystretch}{1.2}

\begin{longtable}{@{}P{1.3cm} P{2.6cm} P{\dimexpr\textwidth-3.9cm-4\tabcolsep\relax}@{}}

\toprule
\textbf{Req.} &
\textbf{Tipo de verificación} &
\textbf{Pruebas y evidencia proporcionada} \\
\midrule
\endfirsthead

\toprule
\textbf{Req.} &
\textbf{Tipo de verificación} &
\textbf{Pruebas y evidencia proporcionada} \\
\midrule
\endhead

RF-01 &
Integración del asistente &
\path{test_wizard_flow.py::test_full_happy_path_produces_valid_params}
recorre los seis pasos y comprueba el diccionario de parámetros resultante;
\path{test_wizard_persisters.py::test_apply_step1_persists_batch_fields}
verifica la persistencia de los parámetros generales;
\path{test_selenium.py::test_step2_configuration_is_propagated_to_step3}
comprueba la propagación de la configuración entre pasos; y
\path{test_wizard_outputs.py::test_back_navigation_preserves_all_choices}
verifica que las opciones se conservan al navegar hacia atrás.
\\[0.9em]

RF-02 &
Pruebas unitarias y de contrato &
\path{test_validators.py::test_step1_rejects_batch_size_that_could_break_generation}
comprueba el límite de 1000 modelos;
\path{test_validators.py::test_step1_rejects_non_numeric_values}
valida los valores numéricos de los parámetros generales;
\path{test_validators.py::test_step2_rejects_minor_levels_without_required_major_levels}
comprueba las dependencias entre niveles;
\path{test_validators.py::test_step3_rejects_invalid_relation_distribution},
\path{test_validators.py::test_step4_rejects_invalid_arithmetic_and_constraint_type_distributions}
y
\path{test_validators.py::test_step5_rejects_invalid_random_attribute_distribution}
verifican las distribuciones de los pasos correspondientes. Las pruebas de
\path{test_params_contract.py} comprueban además el contrato entre los
parámetros del asistente y el modelo de configuración del generador.
\\[0.9em]

RF-03 &
Pruebas unitarias, integración y sistema &
\path{test_engine.py::test_num_models_respected} comprueba el tamaño del lote;
\path{test_engine.py::test_determinism_same_seed_same_output} y
\path{test_engine.py::test_different_seed_different_output} verifican el
comportamiento asociado a la semilla;
\path{test_wizard_outputs.py::test_wrapper_generates_feature_model_with_received_arguments}
comprueba los argumentos entregados a
\texttt{GenerateFeatureModel};
\path{test_wizard_outputs.py::test_wrapper_generate_models_returns_serialized_batch}
valida la generación y serialización de un lote;
\path{test_wizard_outputs.py::test_ensure_satisfiable_retries_until_sat}
comprueba los reintentos de generación satisfacible; y
\path{test_selenium.py::test_step6_generates_one_uvl_model_with_pyodide}
ejecuta el flujo ordinario de generación desde la interfaz.
\\[0.9em]

RF-04 &
Pruebas unitarias e integración &
\path{test_engine.py::test_relation_kinds_create_expected_cardinalities}
comprueba las relaciones \textit{mandatory}, \textit{optional},
\textit{alternative}, \textit{or} y de cardinalidad;
\path{test_engine.py::test_group_cardinality_on_with_weight_produces_group_groups}
y
\path{test_engine.py::test_feature_cardinality_on_produces_some}
validan las cardinalidades;
\path{test_engine.py::test_arithmetic_level_produces_arith_constraints},
\path{test_engine.py::test_aggregate_functions_produce_sum_avg}
y
\path{test_engine.py::test_string_level_produces_len_constraints}
comprueban los niveles aritmético, agregado y de cadenas; y
\path{test_engine.py::test_generated_model_declares_correct_uvl_includes}
verifica las declaraciones \texttt{include} correspondientes.
\\[0.9em]

RF-05 &
Pruebas de sistema e integración &
\path{test_wizard_outputs.py::test_wrapper_generate_one_model_returns_download_payload}
comprueba que el resultado contiene el nombre, el contenido UVL y los
contadores de características y restricciones;
\path{test_selenium.py::test_step6_generates_one_uvl_model_with_pyodide}
verifica que la interfaz muestra el resultado y sus acciones; y
\path{test_selenium.py::test_step6_shows_output_options}
comprueba que las opciones de salida aparecen antes de generar.
\\[0.9em]

RF-06 &
Pruebas unitarias, integración y sistema &
\path{test_validators.py::test_step6_accepts_all_output_options}
comprueba la configuración de los sufijos;
\path{test_wizard_persisters.py::test_apply_step6_output_persists_checkboxes}
verifica su persistencia;
\path{test_wizard_outputs.py::test_wrapper_builds_expected_filename}
y
\path{test_wizard_outputs.py::test_serialize_generated_model_covers_all_filename_branches}
comprueban las combinaciones de nombres y sufijos;
\path{test_wizard_outputs.py::test_wrapper_generate_one_model_returns_download_payload}
valida el contenido preparado para la descarga; y
\path{test_selenium.py::test_step6_generates_one_uvl_model_with_pyodide}
comprueba que la acción de descarga queda disponible tras generar el modelo.
La descarga efectiva del archivo comprimido se verificó manualmente.
\\[0.9em]

\midrule

RNF-01 &
Integración arquitectónica &
\path{test_engine.py::test_generate_feature_model_is_discoverable_by_flamapy}
comprueba que \texttt{DiscoverMetamodels} localiza la operación
\texttt{GenerateFeatureModel};
\path{test_wizard_outputs.py::test_wrapper_build_one_delegates_to_feature_model_generator}
verifica la delegación desde el \textit{wrapper}; y
\path{test_params_contract.py::test_wizard_referenced_params_keys_are_mapped_to_generator_model}
comprueba el contrato entre los parámetros de UVLHub y
\texttt{FmgeneratorModel}.
\\[0.9em]

RNF-02 &
Compatibilidad sintáctica con UVL &
\path{test_wizard_outputs.py::test_o1_level_combinations_roundtrip_to_uvl_and_boolean_sat}
genera combinaciones de niveles, serializa los modelos y los vuelve a leer
mediante \texttt{UVLReader};
\path{test_wizard_outputs.py::test_wrapper_serializes_uvl_and_reads_includes}
comprueba la conservación de los \texttt{include}; y
\path{test_engine.py::test_generated_model_declares_correct_uvl_includes}
verifica que cada configuración declara los niveles UVL necesarios.
\\[0.9em]

RNF-03 &
Ejecución en el navegador e integración con el servidor &
\path{test_selenium.py::test_wizard_reaches_step5_with_pyodide_ready}
comprueba la carga de Pyodide, la instalación de dependencias y la
disponibilidad de la API \texttt{generateOne};
\path{test_selenium.py::test_step6_generates_one_uvl_model_with_pyodide}
ejecuta la generación ordinaria en el navegador;
\path{test_wizard_flow.py::test_generate_sat_returns_models}
comprueba la generación solicitada con satisfacibilidad; y
\path{test_wizard_flow.py::test_params_json_returns_400_without_params}
verifica la respuesta del endpoint de configuración cuando no hay parámetros
en la sesión.
\\[0.9em]

RNF-04 &
Usabilidad, navegación y estado &
\path{test_selenium.py::test_step1_numeric_fields_expose_expected_browser_constraints}
comprueba las restricciones de los campos numéricos;
\path{test_selenium.py::test_step2_type_level_enables_arithmetic_level_in_ui}
y
\path{test_selenium.py::test_step2_arithmetic_toggles_arithmetic_minor_levels_visibility}
verifican el comportamiento de las opciones dependientes;
\path{test_selenium.py::test_wizard_preserves_step1_values_after_next_and_previous}
y
\path{test_wizard_state.py::test_step_state_is_preserved_after_returning_to_previous_step}
comprueban la conservación de los valores; y
\path{test_wizard_flow.py::test_step1_invalid_post_renders_errors}
verifica que los errores de validación se muestran en el formulario.
\\[0.9em]

RNF-05 &
Suite automatizada y cobertura &
La suite automatizada ejecuta 220 pruebas de UVLHub y 44 pruebas del paquete
\texttt{fm\_generator}. Incluye pruebas unitarias en
\path{test_engine.py}, \path{test_validators.py} y
\path{test_params_contract.py}; pruebas de integración en
\path{test_wizard_flow.py}, \path{test_wizard_outputs.py},
\path{test_wizard_persisters.py}, \path{test_wizard_state.py} y
\path{test_wizard_builders.py}; y pruebas de sistema en
\path{test_selenium.py}. La Tabla~\ref{tab:cobertura_ficheros}, presentada en
el Capítulo~\ref{chap:pruebas}, registra una cobertura conjunta del código de
producción del 86,5\,\%, superior al umbral fijado del 80\,\%.
\\

\bottomrule
\caption{Matriz de trazabilidad entre requisitos funcionales y no funcionales
y sus verificaciones.}
\label{tab:trazabilidad_rf_tests}

\end{longtable}
\endgroup

La descarga efectiva del archivo comprimido se comprobó manualmente, ya que
las pruebas automatizadas verifican la preparación de los ficheros y la
disponibilidad de la acción de descarga, pero no inspeccionan el archivo
descargado por el navegador.

Una misma prueba puede aportar evidencia para varios requisitos cuando sus
aserciones comprueban propiedades diferentes del flujo. En cada fila se indica
la propiedad concreta verificada para evitar interpretar esta reutilización
como una duplicación de cobertura. La matriz recoge pruebas representativas y
no pretende enumerar la totalidad de la suite automatizada.


Con los requisitos funcionales, no funcionales, reglas de negocio y restricciones del proyecto definidos y trazados, queda establecido el comportamiento esperado del sistema y las condiciones que debe cumplir la solución desarrollada. El siguiente capítulo aborda cómo satisfacer estos requisitos mediante la definición de la arquitectura del sistema, los componentes principales y las decisiones de diseño adoptadas durante el desarrollo.
